% Part: computability
% Chapter: computability-theory
% Section: complement-ce

\documentclass[../../../include/open-logic-section]{subfiles}

\begin{document}

\olfileid[hi]{cmp}{thy}{cmp}
\olsection{संगणनीयतः परिगणनीय समुच्चय पूरक के अंतर्गत संवृत नहीं हैं}

मान लें कि $A$ संगणनीयतः परिगणनीय है। क्या~$A$ का पूरक
$\Complement{A} = \Nat \setminus A$ भी सदैव संगणनीयतः परिगणनीय होता है?
अगला प्रमेय और उपप्रमेय दिखाते हैं कि उत्तर ``नहीं'' है।

\begin{thm}
\ollabel{thm:ce-comp}
$A$ को प्राकृतिक संख्याओं का कोई समुच्चय मानें। तब $A$ तभी और केवल तभी
संगणनीय है जब $A$ और $\Complement{A}$ दोनों संगणनीयतः परिगणनीय हों।
\end{thm}

\begin{proof}
आगे की दिशा सरल है: यदि $A$ संगणनीय है, तो $\Complement{A}$ भी संगणनीय
है ($\Char{A} = 1 \tsub
\Char{\Complement{A}}$), और इसलिए दोनों
संगणनीयतः परिगणनीय हैं।

दूसरी दिशा में मान लें कि $A$ और~$\Complement{A}$ दोनों संगणनीयतः
परिगणनीय हैं। $A$ को~$\cfind{d}$ का प्रांत और $\Complement{A}$ को
~$\cfind{e}$ का प्रांत मानें। $h$ को परिभाषित करें:
\[
h(x) = \umin{s}{(T(d,x,s) \lor T(e,x,s))}.
\]
दूसरे शब्दों में, आगत~$x$ पर $h$,~$\cfind{d}$ की रुकने वाली संगणना अथवा
~$\cfind{e}$ की रुकने वाली संगणना खोजता है। अब यदि $x \in A$, तो वह
पहली स्थिति में सफल होगा; और यदि $x \in
\Complement{A}$, तो दूसरी स्थिति
में सफल होगा। अतः $h$~एक पूर्ण संगणनीय फलन है। किंतु अब प्रत्येक~$x$ के
लिए $x \in
A$ तभी और केवल तभी जब $T(e, x, h(x))$, अर्थात् जब
$\cfind{e}$ ही परिभाषित हो। चूँकि $T(e, x, h(x))$ एक संगणनीय संबंध है,
$A$~संगणनीय है।
\end{proof}

\begin{explain}
जो हो रहा है उसे अनौपचारिक संगणनात्मक पदों में समझना सरल है: $A$ का
निर्णय करने के लिए आगत $x$ पर $\cfind{e}$ और $\cfind{f}$ की रुकने वाली
संगणनाएँ खोजें। उनमें से एक अवश्य रुकेगी; यदि वह $\cfind{e}$ है, तो $x$,
~$A$ में है, अन्यथा $x$,~$\Complement{A}$ में है।
\end{explain}

\begin{cor}
\ollabel{cor:comp-k}
$\Complement{K_0}$ संगणनीयतः परिगणनीय नहीं है।
\end{cor}

\begin{proof}
हम जानते हैं कि $K_0$ संगणनीयतः परिगणनीय है, किंतु संगणनीय नहीं। यदि
$\Complement{K_0}$ संगणनीयतः परिगणनीय होता, तो \olref{thm:ce-comp} से
$K_0$ संगणनीय होता, जो \olref[ncp]{thm:K-0} के विरुद्ध है।
\end{proof}

\end{document}
